\begin{lemma}[The weak-type level set of a coordinate of $Th$ occupies at most a quarter of a small ball]
Fix $n \geq 1$ and $j \in \{1,\ldots,n\}$, let $\lambda$ be Lebesgue measure on
$\mathbb{R}^n$ and let $\gamma_n$ be the Gaussian measure of
\noderef{GaussMeasure}.  Let
\[
  T : L^1(\gamma_n;\mathbb{R}) \longrightarrow
      \mathcal{M}\bigl(\gamma_n;\mathbb{R}^n\bigr)
\]
be \emph{any} linear map into the $\gamma_n$-almost-everywhere equivalence
classes of measurable $\mathbb{R}^n$-valued functions such that
\begin{enumerate}
\item[(T1)] $T(\iota v) = R_\gamma v$ for every $v \in L^2(\gamma_n)$, with
$\iota$ the inclusion of \noderef{ToL1} and $R_\gamma$ the transform of
\noderef{RieszTransform}; and
\item[(T2)] $t\,\gamma_n\bigl(\{x : |Tf(x)|_{\mathbb{R}^n} > t\}\bigr)
\leq 2\,\|f\|_{L^1(\gamma_n)}$ for every $f \in L^1(\gamma_n;\mathbb{R})$ and
every real $t > 0$.
\end{enumerate}
Let $R \geq 1$.  Then there is a constant $K \in (0,\infty)$, depending only on
$n$, $j$ and $R$, such that the following holds.

Let $h \in L^1(\gamma_n)$, written also for its canonical representative, let
$x \in \mathbb{R}^n$ and let $\varepsilon$ and $\Phi$ be real, and suppose
\[
  |x| \leq R, \qquad 0 < \varepsilon \leq 1, \qquad 0 \leq \Phi ,
  \qquad
  \int_{\mathbb{R}^n}|h|\,d\gamma_n \;\leq\; \Phi\cdot\lambda\bigl(B(x,\varepsilon)\bigr) .
\]
Then
\[
  \gamma_n\Bigl(\bigl\{z \in B(x,\tfrac12\varepsilon)
      : \bigl|(Th)_j(z)\bigr| > K\Phi\bigr\}\Bigr)
    \;\leq\; \tfrac14\,\gamma_n\bigl(B(x,\tfrac12\varepsilon)\bigr) ,
\]
where $(Th)_j$ denotes the $j$th coordinate of the canonical representative of
$Th$, and $\lambda(B(x,\varepsilon))$ is read as the real number it is, the
ball being bounded and hence of finite measure.

\emph{The threshold is $K\Phi$ and not a named constant, deliberately.}  A
consumer arrives with a specific threshold built from its own bookkeeping ---
in the intended application, $8\cdot2^n\pi^{-n/2}e^{(R+2)^2}\Phi$ --- and the
only thing that matters about it is that it is a large enough multiple of
$\Phi$.  Exporting the multiple existentially rather than writing the
consumer's expression into this statement keeps the node free of arithmetic
that is none of its business, and the existential constant absorbs it at no
cost.  One may take $K = 8\cdot2^n/a$ with
$a = \pi^{-n/2}e^{-(R+2)^2}$.

\emph{The degenerate case $\Phi = 0$ is included and is not vacuous.}  When
$\Phi = 0$ the hypothesis forces $\int|h|\,d\gamma_n = 0$, hence $h = 0$ in
$L^1(\gamma_n)$, and the conclusion asserts that
$\{z \in B(x,\tfrac12\varepsilon) : |(Th)_j(z)| > 0\}$ is $\gamma_n$-null.
That is true but it needs the \emph{linearity} of $T$ --- it is $T0 = 0$ that
makes it true --- and not (T2), which is unavailable at threshold $0$.  So the
statement is not merely the weak-type bound rearranged: it covers a case the
weak-type bound cannot reach.

\emph{Both the hypothesis and the conclusion are insensitive to the
representative.}  Changing the representative of $h$ on a $\gamma_n$-null set
changes neither $\int|h|\,d\gamma_n$ nor the class $Th$, and changing the
representative of $Th$ on a null set changes neither the level set's
$\gamma_n$-measure nor the conclusion; so fixing the canonical representatives
is notation, not an extra hypothesis.

\emph{The constant depends on $R$, necessarily, and no uniformity is claimed.}
$K$ degrades in $R$ through $a$, the lower bound for the Gaussian density on
$\overline B(0,R+2)$, and in $n$ through both $a$ and the factor $2^n$ coming
from the ratio $\lambda(B(x,\tfrac12\varepsilon))/\lambda(B(x,\varepsilon))$.
\end{lemma}
\begin{proof}
Put $a = \pi^{-n/2}e^{-(R+2)^2} > 0$ and $K = 8\cdot2^n/a \in (0,\infty)$,
which depends only on $n$, $j$ and $R$.  Fix $h$, $x$, $\varepsilon$, $\Phi$
and the hypotheses, and write $V = \lambda(B(x,\varepsilon)) \in [0,\infty)$,
finite because the ball is bounded, and
$B' = B(x,\tfrac12\varepsilon)$.

\emph{Step 1: three facts about the balls.}  Lebesgue measure on
$\mathbb{R}^n$ is an additive Haar measure and $\mathbb{R}^n$ has real
dimension $n$, so $\lambda(B(x,r)) = r^n\lambda(B(0,1))$ for every $r \geq 0$
(shared context; the hypothesis $n \geq 1$ is what makes $\mathbb{R}^n$
nontrivial, as that statement requires).  Taking $r = \tfrac12\varepsilon$ and
$r = \varepsilon$ and dividing,
\begin{equation*}
  \lambda(B') = 2^{-n}\,V . \tag{1}
\end{equation*}
Next, $B'$ is measurable and $B' \subseteq \overline B(0,R+2)$: for $y \in B'$
we have $|y| - |x| \leq |y-x| < \tfrac12\varepsilon \leq \tfrac12$, so
$|y| \leq R + \tfrac12 \leq R+2$.  Hence clause (2) of
\noderef{GaussLebesgueBallComparison}, applied with $R+2$ in place of its $R$
and with $E = B'$, gives
\begin{equation*}
  a\,\lambda(B') \leq \gamma_n(B') . \tag{2}
\end{equation*}
Combining $(1)$ and $(2)$ and dividing by $4$,
\begin{equation*}
  a\,2^{-n}\,V/4 \;\leq\; \tfrac14\gamma_n(B') . \tag{3}
\end{equation*}
Finally, the $L^1(\gamma_n)$ norm of $h$ is the integral of the absolute value
of its canonical representative, so the hypothesis reads
\begin{equation*}
  \|h\|_{L^1(\gamma_n)} \leq \Phi\,V . \tag{4}
\end{equation*}

\emph{Step 2: the case $\Phi = 0$.}  Then $(4)$ gives
$\|h\|_{L^1(\gamma_n)} \leq 0$, and a norm is nonnegative, so
$\|h\|_{L^1(\gamma_n)} = 0$ and therefore $h = 0$ in $L^1(\gamma_n)$.  Since
$T$ is linear, $Th = T0 = 0$, so the canonical representative of $Th$ is $0$
$\gamma_n$-almost everywhere.  The threshold is $K\Phi = 0$, so the level set
in question is $\{z \in B' : |(Th)_j(z)| > 0\}$, and every $z$ in it has
$(Th)_j(z) \neq 0$, hence $(Th)(z) \neq 0$.  So the level set is contained in
the $\gamma_n$-null set $\{z : (Th)(z) \neq 0\}$ and is therefore itself
$\gamma_n$-null, which is at most $\tfrac14\gamma_n(B')$.

\emph{Step 3: the case $\Phi > 0$.}  Then $t := K\Phi > 0$.  Write
\[
  E_1 = \{z \in B' : |(Th)_j(z)| > t\}, \qquad
  F = \{z : |(Th)_j(z)| > t\} ,
\]
so $E_1 \subseteq F$ and hence $\gamma_n(E_1) \leq \gamma_n(F)$ by
monotonicity of the measure.  Applying \noderef{CoordinateEndpointBound} to
this $T$ with its two properties, which are (T1) and (T2) and hold by
assumption, to the element $h$, to the coordinate $j$ and to this $t > 0$,
\[
  t\,\gamma_n(F) \leq 2\,\|h\|_{L^1(\gamma_n)} .
\]
Since $t > 0$ is finite and nonzero, multiplying through by $t^{-1}$ gives
$\gamma_n(F) \leq 2\|h\|_{L^1(\gamma_n)}/t$.  By $(4)$ and then the definition
of $K$,
\[
  \frac{2\|h\|_{L^1(\gamma_n)}}{t}
    \leq \frac{2\,\Phi V}{K\Phi}
    = \frac{2V}{K}
    = \frac{2V\,a}{8\cdot2^n}
    = \frac{a\,2^{-n}V}{4} ,
\]
the cancellation of $\Phi$ being legitimate because $\Phi > 0$.  Chaining this
with $\gamma_n(E_1) \leq \gamma_n(F)$ and with $(3)$,
\[
  \gamma_n(E_1) \leq \frac{a\,2^{-n}V}{4} \leq \tfrac14\gamma_n(B') ,
\]
which is the assertion.
\end{proof}
